% Section 6.4, from lectures/L06/appendix/c_exercises.md.

\section{Exercises}
\label{c6:sec:exercises}\label{c6:app:c}
These exercises reinforce the concepts from \secref{c6:app:a} and \secref{c6:app:b}.

\begin{aside}
\textbf{How to check your work.} Write each exercise set as its own application under
\file{lectures/L06/exercises}, with the file names given there, and run \code{make test} from the
repository root. A set’s tests switch on as soon as its files exist, and a set you have not started
is reported as \code{SKIP}. Each application grows through its set, and the tests check where it
ends up. Shell commands are fed to your program on its standard input, and its output is checked for
the lines each exercise shows, in order; the prompt, the echoed commands and anything else in
between do not matter.

The solutions are published in the companion repository, under
\file{lectures/L06/appendix/solutions}, a few days after the lecture. Try each exercise yourself
before you read them.
\end{aside}

Remember \code{source env.sh} in every new terminal. Build and run an exercise from the repository
root:

\begin{shell}
west build -p -b native_sim/native/64 lectures/L06/exercises/exercise1
west build -t run
\end{shell}

% ---------------------------------------------------------------------
\exerciseset{Shell commands of your own}
\label{c6:set:1}
Create an application named \code{exercise1} to work in for this exercise set.

\exercise{The setpoint}{Code}
\label{c6:ex:1.1}
In this exercise you will give a thermostat a shell command that reads and changes its setpoint.

\task{Tasks}
\begin{itemize}
  \item In \code{prj.conf}, enable the shell, and put the shell’s UART on the terminal the program is
    started from.
  \item In \code{src/main.c}, keep the setpoint in a \code{static int setpoint_c}, initially 21.
  \item Create a root command \code{setpoint} with the help text \code{The thermostat's setpoint.}, no
    handler of its own, and two subcommands:
\begin{itemize}
  \item \code{get}, which takes no arguments, and prints \code{Setpoint: <t> C}.
  \item \code{set <celsius>}, which takes exactly one argument, sets the setpoint, and prints \code{Setpoint
    set to <t> C}.
\end{itemize}
  \item Let the shell check the argument counts: declare them with \code{SHELL_CMD_ARG()}.
  \item \code{main()} prints \code{Thermostat ready} and returns.
\end{itemize}

Pipe a session into the program:

\begin{shell}
printf 'setpoint get\nsetpoint set 25\nsetpoint get\n' \
    | ./build/zephyr/zephyr.exe -stop_at=1
\end{shell}

\noindent The output must contain, in this order:

\begin{console}
Thermostat ready
Setpoint: 21 C
Setpoint set to 25 C
Setpoint: 25 C
\end{console}

\task{Reflection}
\begin{itemize}
  \item The program’s output also contains \code{uart:~$ setpoint get}. Where does that line come from?
  \item Run \code{setpoint} alone, and then \code{retval}. What does \code{retval} print, and why that
    value?
  \item \code{main()} returns at once. Why does the shell keep answering commands?
\end{itemize}

\exercise{Wrong input}{Code}
\label{c6:ex:1.2}

\task{Tasks}
Make \code{setpoint set} reject what is not a valid setpoint:

\begin{itemize}
  \item A word that is not a whole decimal number, such as \code{warm} or \code{25x}, prints the error
    \code{not a temperature: '<word>'} with \code{shell_error()} and returns \code{-EINVAL}.
  \item A number outside 5 to 30 prints the error \code{setpoint must be between 5 and 30 C, not <t>} and
    returns \code{-ERANGE}. The setpoint does not change.
\end{itemize}

For the session

\begin{console}
setpoint get
setpoint set 25
setpoint get
setpoint set 31
retval
setpoint set warm
retval
setpoint set 25x
setpoint set
\end{console}

the output must contain, in this order:

\begin{console}
Setpoint: 21 C
Setpoint set to 25 C
Setpoint: 25 C
setpoint must be between 5 and 30 C, not 31
-34
not a temperature: 'warm'
-22
not a temperature: '25x'
set: wrong parameter count
\end{console}

\noindent \textbf{Hint:} \code{strtol()} from \code{<stdlib.h>} sets its \code{end} argument to the
first character it did not convert.

\task{Reflection}
\begin{itemize}
  \item Why is \code{atoi()} not good enough here?
  \item Which part of the output does the shell print by itself, without your handler being called?
  \item \code{setpoint set 99999999999999999999} is a number too. What does your handler do with it, and
    why?
\end{itemize}

% ---------------------------------------------------------------------
\exerciseset{A setpoint that survives a restart}
\label{c6:set:2}
Copy \code{exercise1} to a new application named \code{exercise2} to work in for this exercise set.

\exercise{Storing it}{Code}
\label{c6:ex:2.1}
In this exercise you will store the setpoint with the settings subsystem, on the simulated flash.

\task{Tasks}
\begin{itemize}
  \item In \code{prj.conf}, enable the settings subsystem with the ZMS backend, and what ZMS needs: flash
    and the flash map.
  \item Keep the setpoint in an \code{int32_t}, stored under the name \code{thermostat/setpoint}.
  \item Register a settings handler for the subtree \code{thermostat} with
    \code{SETTINGS_STATIC_HANDLER_DEFINE()}. Its \code{h_set} callback takes the stored \code{setpoint},
    and rejects a stored value whose length is not that of an \code{int32_t} with \code{-EINVAL}.
  \item In \code{main()}, initialize the settings subsystem, load the settings, and print \code{Loaded
    setpoint: <t> C} instead of \code{Thermostat ready}.
  \item \code{setpoint set} stores the new value with \code{settings_save_one()} before it changes the
    setpoint in RAM. If storing fails, it prints \code{cannot store the setpoint: <error>} with
    \code{shell_error()} and returns the error.
\end{itemize}

Run the program twice, from the same directory:

\begin{shell}
cd build
printf 'setpoint set 25\n' | ./zephyr/zephyr.exe -stop_at=1
printf 'setpoint get\n' | ./zephyr/zephyr.exe -stop_at=1
\end{shell}

\noindent The first run must print \code{Loaded setpoint: 21 C} and \code{Setpoint set to 25 C}, and
the second \code{Loaded setpoint: 25 C} and \code{Setpoint: 25 C}.

\task{Reflection}
\begin{itemize}
  \item Where is the setpoint stored? Find the file, and find the setting’s name in it with \code{strings}.
  \item Run the second command again from the repository root instead of from \code{build/}. What does it
    load, and why?
  \item Run it with \code{-flash_in_ram}. What does it load, and has the file changed?
  \item In \code{h_set}, why read the value into a local variable and check it, rather than straight into
    the setpoint?
\end{itemize}

\exercise{Forgetting it}{Code}
\label{c6:ex:2.2}

\task{Tasks}
Add a subcommand \code{setpoint reset}, which takes no arguments, deletes the stored setpoint with
\code{settings_delete()}, sets the setpoint back to 21, and prints \code{Setpoint reset to 21 C}.

In five runs from the same directory, where each run is given the commands shown, the output must
contain:

\begin{booktable}{@{}lLL@{}}
  \thead{Run} & \thead{Commands} & \thead{Output} \\
  \midrule
  1 & \code{setpoint set 25} & \code{Loaded setpoint: 21 C}, then \code{Setpoint set to 25 C} \\
  2 & \code{setpoint get} & \code{Loaded setpoint: 25 C}, then \code{Setpoint: 25 C} \\
  3 & none, run with \code{-flash_in_ram} & \code{Loaded setpoint: 21 C} \\
  4 & \code{setpoint reset} & \code{Loaded setpoint: 25 C}, then \code{Setpoint reset to 21 C} \\
  5 & none & \code{Loaded setpoint: 21 C} \\
\end{booktable}

\task{Reflection}
\begin{itemize}
  \item Enable \code{CONFIG_SETTINGS_SHELL} on the command line, and run \code{settings list} and
    \code{settings read thermostat/setpoint} after setting the setpoint to 27. What do the four bytes it
    prints mean?
  \item How would you make every run of the tests start from a factory-new device, without deleting files by
    hand?
\end{itemize}

% ---------------------------------------------------------------------
\exerciseset{Fatal errors}
\label{c6:set:3}
Create an application named \code{exercise3} to work in for this exercise set.

\exercise{An impossible reading}{Code}
\label{c6:ex:3.1}
In this exercise you will make a thread give up with \code{k_oops()}, and see what that costs
without a handler of your own.

\task{Tasks}
\begin{itemize}
  \item In \code{prj.conf}, enable logging and thread names.
  \item In \code{src/main.c}, register the log module \code{app}, and define the readings:
\begin{ccode}
static const int readings[] = {21, 22, -300, 23};
\end{ccode}
  \item Write a thread function \code{sampler()} that goes through the readings. A reading below -273 logs
    the error \code{Impossible reading: <t> C} and calls \code{k_oops()}; any other logs \code{Sample
    <t> C}.
  \item In \code{main()}, create the thread with \code{k_thread_create()}, with a 1024-byte stack from
    \code{K_THREAD_STACK_DEFINE()} and priority 5, name it \code{sampler}, and wait for it with
    \code{k_thread_join()}. Then log \code{main() carries on}, and return.
\end{itemize}

Build, run, and check the exit status:

\begin{shell}
./build/zephyr/zephyr.exe -stop_at=1; echo "exit status $?"
\end{shell}

\task{Reflection}
\begin{itemize}
  \item Which lines does the kernel print, and what does each one tell you?
  \item Does \code{main() carries on} appear? What is the exit status, and why?
  \item Build once more with \code{CONFIG_LOG=n} on the command line. What does a fatal error print now?
\end{itemize}

\exercise{A handler of your own}{Code}
\label{c6:ex:3.2}

\task{Tasks}
\begin{itemize}
  \item Create \code{src/fatal.c}, with a log module named \code{fatal}, and define
    \code{k_sys_fatal_error_handler()} in it. It shall:
\begin{itemize}
  \item Switch logging to synchronous with \code{LOG_PANIC()}.
  \item Log the error \code{Fatal error <reason> in thread '<name>'}.
  \item For an oops, log \code{Aborting the thread} and return.
  \item For anything else, log \code{Halting the system} and call \code{k_fatal_halt()}.
\end{itemize}
  \item Add \code{src/fatal.c} to the application in \code{CMakeLists.txt}.
\end{itemize}

The output must now contain, in this order, and the exit status must be 0:

\begin{console}
<inf> app: Sample 21 C
<inf> app: Sample 22 C
<err> app: Impossible reading: -300 C
>>> ZEPHYR FATAL ERROR 3: Kernel oops
<err> fatal: Fatal error 3 in thread 'sampler'
<err> fatal: Aborting the thread
<inf> app: main() carries on
\end{console}

\task{Reflection}
\begin{itemize}
  \item Why does \code{Sample 23 C} never appear?
  \item Why does \code{k_thread_join()} return, when the thread never returned from its function?
  \item Remove the check for the oops, so that the handler always returns. Nothing changes in this program’s
    output. Why not, and why is it still wrong?
\end{itemize}

\exercise{An assertion}{Code}
\label{c6:ex:3.3}

\task{Tasks}
\begin{itemize}
  \item In \code{src/main.c}, add a function:
\begin{ccode}
static void set_heater_duty(int percent)
{
  __ASSERT((percent >= 0) && (percent <= 100), "duty cycle %d %% is out of range", percent);
  LOG_INF("Heater duty cycle %d %%", percent);
}
\end{ccode}
  \item At the end of \code{main()}, after \code{main() carries on}, call \code{set_heater_duty(150)}.
\end{itemize}

Built as usual, the program must also print \code{<inf> app: Heater duty cycle 150 \%} at the end,
and exit with status 0. Built with \code{-DCONFIG_ASSERT=y}, the output must contain, in this order,
and the exit status must be 1:

\begin{console}
<inf> app: main() carries on
duty cycle 150 % is out of range
>>> ZEPHYR FATAL ERROR 4: Kernel panic
<err> fatal: Fatal error 4 in thread 'main'
<err> fatal: Halting the system
\end{console}

\task{Reflection}
\begin{itemize}
  \item Why does the same program pass a duty cycle of 150 \% in one build and halt in the other?
  \item With the handler from the last reflection of Exercise~\ref{c6:ex:3.2}, the one that always returns,
    build with \code{-DCONFIG_ASSERT=y} and run. What happens, and what is the exit status?
  \item Why would \code{__ASSERT(settings_save_one(...) == 0, "...")} be a bug?
\end{itemize}

% ---------------------------------------------------------------------
\exerciseset{Sanitizers}
\label{c6:set:4}
Create an application named \code{exercise4} to work in for this exercise set.

\exercise{A battery monitor}{Code}
\label{c6:ex:4.1}
In this exercise you will find two bugs with the sanitizers, and fix them.

\task{Tasks}
Create \code{CMakeLists.txt}, a \code{prj.conf} with only a comment, and a \code{src/main.c} holding
exactly this:

\begin{ccode}
/*
 * SPDX-License-Identifier: MIT
 *
 * Exercise Set 4: a battery monitor, with two bugs.
 */

#include <zephyr/kernel.h>

/* How many of the most recent readings the average is taken over. */
#define WINDOW 4

/* The average of the first n readings in a window. */
static __noinline int32_t average(const int32_t *window, size_t n)
{
	int32_t sum = 0;

	for (size_t i = 0; i <= n; i++) {
		sum += window[i];
	}
	return sum / (int32_t)n;
}

/* The power drawn at a voltage and a current, in microwatts. */
static __noinline int32_t power_uw(int32_t millivolts, int32_t milliamps)
{
	return millivolts * milliamps * 1000;
}

static const int32_t readings_mv[] = {3300, 3310, 3290, 3300, 5000, 5020};
static const int32_t load_ma = 450;

int main(void)
{
	int32_t window[WINDOW] = {0};

	for (size_t i = 0; i < ARRAY_SIZE(readings_mv); i++) {
		window[i % WINDOW] = readings_mv[i];

		printk("%d mV, average %d mV, %d uW\n", readings_mv[i],
		       average(window, MIN(i + 1, WINDOW)), power_uw(readings_mv[i], load_ma));
	}
	return 0;
}
\end{ccode}

\noindent Build and run it as it is.

\task{Reflection}
\begin{itemize}
  \item Which lines of the output are wrong? Would you have noticed if you did not know what to expect?
  \item The first three lines are right, although the bug in \code{average()} is there from the first call.
    Why do they come out right?
\end{itemize}

\exercise{AddressSanitizer}{Code}
\label{c6:ex:4.2}

\task{Tasks}
Build with \code{-DCONFIG_ASAN=y}, run, and fix the bug that AddressSanitizer reports, changing as
little as possible.

\task{Reflection}
\begin{itemize}
  \item Which line, which variable, and which kind of access does the report name?
  \item What is the exit status, and why do the lines after the report never appear?
\end{itemize}

\exercise{UndefinedBehaviorSanitizer}{Code}
\label{c6:ex:4.3}

\task{Tasks}
Build with \code{-DCONFIG_UBSAN=y}, run, and fix the bug that UndefinedBehaviorSanitizer reports, so
that \code{power_uw()} returns an \code{int64_t} and computes the power without overflow for any
voltage and current that fit in an \code{int32_t} each.

Built with \code{-DCONFIG_ASAN=y -DCONFIG_UBSAN=y}, and built without either, the program must print
exactly:

\begin{console}
3300 mV, average 3300 mV, 1485000000 uW
3310 mV, average 3305 mV, 1489500000 uW
3290 mV, average 3300 mV, 1480500000 uW
3300 mV, average 3300 mV, 1485000000 uW
5000 mV, average 3725 mV, 2250000000 uW
5020 mV, average 4152 mV, 2259000000 uW
\end{console}

\task{Reflection}
\begin{itemize}
  \item The program did not stop at the report. How is that different from AddressSanitizer?
  \item Would \code{return (int64_t)(millivolts * milliamps * 1000);} fix the bug? Try it with UBSan.
  \item Build the fixed program with \code{-DCONFIG_COVERAGE=y}, run it, and run \code{gcov} on
    \code{main.c} as in \secref{c6:app:b}. How many times did the body of the loop in \code{average()}
    run, and why that number?
\end{itemize}

% ---------------------------------------------------------------------
\exerciseset{Sockets}
\label{c6:set:5}
Create an application named \code{exercise5} to work in for this exercise set.

\exercise{A UDP server}{Code}
\label{c6:ex:5.1}
In this exercise you will write a UDP server that answers every datagram in capitals, and a client
for it in the same program, talking through the host’s network with offloaded sockets.

\task{Tasks}
\begin{itemize}
  \item In \code{prj.conf}, enable networking with IPv4 only, sockets, and \code{native_sim}’s offloaded
    sockets, as in the \file{lectures/L06/examples/sockets} example. No logging is needed.
  \item Define \code{PORT} as 4300.
  \item Write a thread \code{server}, with \code{K_THREAD_DEFINE()}, a 2048-byte stack and priority 7,
    which:
\begin{itemize}
  \item Creates a UDP socket, and binds it to \code{PORT} on every address (\code{NET_INADDR_ANY}).
  \item In an endless loop, receives a datagram of up to 64 bytes with \code{zsock_recvfrom()}, turns it
    into capitals with \code{toupper()}, and sends it back with \code{zsock_sendto()} to the address it
    came from.
\end{itemize}
  \item Use the prefixed names throughout: \code{struct net_sockaddr_in}, \code{NET_AF_INET},
    \code{net_htons()}, and so on.
\end{itemize}

\exercise{The client}{Code}
\label{c6:ex:5.2}

\task{Tasks}
In \code{main()}:

\begin{itemize}
  \item Create a UDP socket, and set its receive timeout to one second with \code{ZSOCK_SO_RCVTIMEO}.
  \item Sleep 100 ms, so that the server, which has a lower priority, has bound its port.
  \item For each of the words \code{hello}, \code{zephyr} and \code{native_sim}, send the word to
    \code{127.0.0.1} at \code{PORT}, receive the reply, and print \code{Sent '<word>', received
    '<reply>'}. If no reply arrives, print \code{No reply to '<word>': <errno>} instead.
  \item Close the socket and print \code{Done}.
\end{itemize}

The output must contain, in this order:

\begin{console}
Sent 'hello', received 'HELLO'
Sent 'zephyr', received 'ZEPHYR'
Sent 'native_sim', received 'NATIVE_SIM'
Done
\end{console}

\task{Reflection}
\begin{itemize}
  \item Run the program without \code{-no-rt} and with \code{-stop_at=10}, and send it a datagram from
    another terminal with \code{printf 'from the host' | nc -u -w1 127.0.0.1 4300}. What comes back, and
    what does that say about where the port is?
  \item Why does the client set a receive timeout, when the example’s TCP client does not?
  \item The output also begins with \code{WARNING: Using a test - not safe - entropy source}. Where does
    that come from?
\end{itemize}
